\section{Resumen ejecutivo}

El 19 de octubre de 2026 se evaluó la seguridad del prototipo \textbf{SOS Clinic},
un sistema de apoyo al registro de pacientes y al proceso de triage compuesto por
un frontend estático, seis microservicios Java, un gateway, descubrimiento con
Eureka y una base PostgreSQL local. La evaluación se ejecutó exclusivamente sobre el entorno local del
repositorio, con datos sintéticos y correo capturado en Mailpit, siguiendo las
fases de PTES y los marcos OWASP WSTG, OWASP API Security Top 10 2023, OWASP
ASVS 4.0, ISO/IEC 27002:2022, CVSS 4.0 y STRIDE/DREAD.

El resultado de la evaluación base es contundente: \textbf{el riesgo global era
crítico}. En la versión evaluada, la aplicación no autenticaba las APIs de
pacientes y triage, no autorizaba por rol las rutas administrativas y conservaba
un secreto JWT por defecto en el código. Esas tres condiciones, combinadas,
permitían que una persona sin credenciales leyera y modificara información clínica,
y que cualquiera con acceso al repositorio emitiera un token de administrador. La
evaluación registró 30 hallazgos brutos que, tras consolidar duplicados, equivalen
a \textbf{25 temas únicos}.

\begin{table}[htbp]
\centering
\caption{Cifras principales de la evaluación.}
\label{tab:cifras}
\begin{tabular}{@{}L{8.2cm}C{3.2cm}C{3.2cm}@{}}
\toprule
\textbf{Indicador} & \textbf{Valor} & \textbf{Observación}\\
\midrule
Hallazgos consolidados & 25 & 30 registros brutos\\
Vulnerabilidades únicas en dependencias & 22 & 4 críticas, 12 altas, 6 medias\\
Capturas de evidencia & 16 & Solicitudes, respuestas y escaneos\\
Servicios y componentes evaluados & 10 & 6 microservicios, frontend, gateway e infraestructura\\
Fases ejecutadas & 6 & Fases 0 a 5 del plan\\
\bottomrule
\end{tabular}
\end{table}

\begin{table}[htbp]
\centering
\caption{Distribución de los 25 hallazgos consolidados en la evaluación base, por severidad.}
\label{tab:severidad}
\begin{tabular}{@{}L{4.2cm}C{2.4cm}L{7.6cm}@{}}
\toprule
\textbf{Severidad} & \textbf{Cantidad} & \textbf{Descripción general}\\
\midrule
\crit{Crítica} & 6 & Compromiso directo de confidencialidad, integridad o autenticación.\\
\alta{Alta} & 8 & Explotación viable con impacto significativo sobre datos o sesiones.\\
\media{Media} & 7 & Debilidades de configuración, exposición o abuso acotado.\\
\baja{Baja} & 2 & Problemas de robustez o de defensa en profundidad.\\
\obs{Observación} & 2 & Desajustes funcionales o falsos positivos documentados.\\
\midrule
\textbf{Total} & \textbf{25} & \\
\bottomrule
\end{tabular}
\end{table}

\needspace{12\baselineskip}
\subsection{Hallazgos principales}

\begin{table}[htbp]
\centering
\caption{Hallazgos críticos que determinan el riesgo global.}
\label{tab:criticos-resumen}
\begin{tabular}{@{}L{1.3cm}L{7.5cm}L{6.0cm}@{}}
\toprule
\textbf{ID} & \textbf{Hallazgo} & \textbf{Impacto principal}\\
\midrule
H-01 & Secreto JWT por defecto y forja de tokens & Emisión de tokens administrativos sin credenciales.\\
H-02 & Lectura de pacientes y triage sin autenticación & Exposición de datos clínicos y de la lista de espera.\\
H-03 & Escritura y cambios de estado sin autenticación & Alteración de historias, episodios y altas.\\
H-04 & Escalada de privilegios en \texttt{/admin/**} & Un usuario de triage crea cuentas administradoras.\\
H-05 & Archivo \texttt{env} con credenciales reales & Exposición de Gmail y Groq; requiere rotación.\\
H-06 & Dependencias críticas vulnerables & Tomcat y FreeMarker con CVEs críticas conocidas.\\
\bottomrule
\end{tabular}
\end{table}

\subsection{Cadena de ataque observada}

La combinación de los hallazgos permite construir una cadena de ataque sencilla
y de alto impacto:

\begin{enumerate}
  \item \textbf{Acceso al repositorio.} El secreto JWT por defecto está versionado
  en \path{application.properties}. Con ese valor se firma un token HS256 con
  perfil de administrador y se accede a \path{/ms-seguridad/admin/**} y a la
  analítica (H-01).
  \item \textbf{Acceso sin credenciales.} Las APIs de pacientes y triage no
  instalan filtro de autenticación. Sin token se leen pacientes, episodios,
  signos vitales e historias PDF, y se crean o modifican registros (H-02, H-03,
  H-09).
  \item \textbf{Robo de sesión.} El frontend inserta nombres y descripciones con
  \texttt{innerHTML}. Un XSS almacenado puede leer el JWT de
  \texttt{sessionStorage} y reutilizarlo (H-07).
  \item \textbf{Abuso del agente IA.} El agente responde sin token, consulta las
  APIs clínicas y consume la cuota de Groq, lo que añade fuga de datos y costo
  operativo (H-10, H-21).
\end{enumerate}

\nota{\textbf{Mensaje clave.} En la versión evaluada no se trataba de una suma de
fallos aislados: la ausencia de autenticación en los servicios de negocio, la
ausencia de autorización por rol en la administración y un secreto compartido con
valor por defecto convertían al prototipo en un sistema \textit{abierto por diseño}
para cualquiera que alcanzara el gateway o el repositorio.}

\subsection{Estado posterior a la remediación}

Después de la evaluación, en la rama \texttt{main} se aplicaron las correcciones
documentadas en \path{seguridad/REMEDIACION-2026-10-19.md}. El resultado es
\textbf{22 hallazgos remediados, 1 mitigado (H-05), 1 parcial (H-18) y 1 documentado
(H-25)}. El riesgo residual disminuye de forma significativa, pero persisten
limitaciones: el logout sigue siendo local, el rate limiting está en memoria, H-09
se cerró por rol y no por propiedad de cada episodio, la CSP conserva
\texttt{unsafe-inline}, el entorno local usa HTTP sin HSTS y los secretos antiguos
siguen en el historial de Git. El detalle está en la sección~\ref{sec:estado-remediacion}.

\subsection{Recomendaciones inmediatas}

\begin{enumerate}
  \item Convertir \texttt{ms-pacientes} y \texttt{ms-triage} en \textit{resource
  servers} con el mismo filtro JWT, y validar el token de forma centralizada en
  el gateway. Aplicar denegación por defecto.
  \item Restringir \path{/ms-seguridad/admin/**} con
  \texttt{hasRole('ADMINISTRADOR')} y separar la gestión de usuarios del resto de
  la administración.
  \item Eliminar el valor por defecto de \path{app.jwt.secret}, exigir
  \texttt{APP\_JWT\_SECRET} desde el entorno y rotar el secreto.
  \item Mover el archivo \texttt{env} fuera del repositorio y rotar la contraseña
  de aplicación de Gmail y la clave de Groq.
  \item Actualizar Tomcat a 11.0.25 o superior, FreeMarker a 2.3.35, Jackson y
  HttpComponents a las versiones corregidas; repetir el SCA.
  \item Sustituir \texttt{innerHTML} por \texttt{textContent} o
  \texttt{createElement}, añadir una CSP sin \texttt{unsafe-inline} y aplicar
  rate limiting en login y 2FA.
\end{enumerate}



